\documentclass[12pt]{article}
\usepackage{graphicx} % Required for inserting images
\usepackage[margin=0.5in]{geometry}
\usepackage{amsmath, amssymb}
\usepackage{mathrsfs}


\usepackage{soul}


% Dr. David Brown Commands
\newcommand{\ds}{\displaystyle}
\newcommand{\R}{\mathbb{R}}
\newcommand{\N}{\mathbb{N}}
\newcommand{\Q}{\mathbb{Q}}
\newcommand{\C}{\mathbb{C}}


% Commands to get script r
\usepackage{calligra}
\DeclareMathAlphabet{\mathcalligra}{T1}{calligra}{m}{n}
\DeclareFontShape{T1}{calligra}{m}{n}{<->s*[2.2]callig15}{}
\newcommand{\scripty}[1]{\ensuremath{\mathcalligra{#1}}}

% Unit commands
\newcommand{\degree}{\text{°}}
\newcommand{\unitSpeed}{\frac{\text{m}}{\text{s}}}
\newcommand{\unitAcceleration}{\frac{\text{m}}{\text{s}^2}}

% Constant commands
\newcommand{\avogadro}{6.022\times10^{23}}

\newcommand{\boltzmannConstK}{1.381\times10^{-23}\frac{\text{J}}{\text{K}}}
\newcommand{\boltzmannConstInEV}{8.617\times10^{-5}\frac{\text{eV}}{\text{K}}}
\newcommand{\planckConst}{6.626\times10^{-34}\text{J}\cdot\text{s}}

\newcommand{\atmP}{1.01325\times10^5\,\text{Pa}}

% Commands to easily write partial derivatives
\newcommand{\partDeriv}[2]{\frac{\partial #1}{\partial #2}}
\newcommand{\doublePartDeriv}[2]{\frac{\partial^2 #1}{\partial^2 #2}}


\title{Problem 6.23}
\author{Gabriel Decker}


\begin{document}



\maketitle
\begin{flushleft}

In this problem, we deal with the rotational motion of diatomic molecules.
In this system, molecules can have rotational energies of $E(j)=j(j+1)\epsilon$ which have multiplicity $2j+1$, where $j\in\N$.
Therefore, the partition function $Z$ is given by the following equation.
\begin{equation}
    Z_\text{rot}=\sum_{j=0}^\infty(2j+1)e^{-j(j+1)\epsilon/kt}
\end{equation}

Using the given $\epsilon=0.00024\,\text{eV}$ for CO and $T=300\,\text{K}$, we will calculate $Z$ in two different ways.
First, I will calculate it using the above equation, using the first many terms (using python).
Then, I will use the high-temperature limit-assumed approximation in the following equation.
\begin{equation}
    Z_\text{rot}\approx\frac{kT}{\epsilon}\ \ \ (kT\gg\epsilon)
\end{equation}\\~\\

First, the brute-force evaluation of the first several terms of the infinite summation was accomplished using the python script found in this directory. I summed the terms between $j=0$ and $j=40$. The following was the result.
$$Z\approx\sum_{j=0}^{40}(2j+1)e^{-j(j+1)\epsilon/kt}=108.046$$
I also tried with the maximum $j$ being 150, and that yielded the same $Z$ up to the three displayed decimal places I displayed.\\~\\

Let's now compare the brute-force result with the high-temperature approximation.
$$Z_\text{rot}\approx\frac{kT}{\epsilon}\ \ \ (kT\gg\epsilon)$$
$$\implies Z_\text{rot}\approx\frac{(\boltzmannConstInEV)(300\,\text{K})}{0.00024\,\text{eV}}$$
$$\implies Z_\text{rot}\approx107.713$$\\~\\

These are clearly not the same answer, but they are pretty close.



\end{flushleft}

\end{document}
